\section{带权矩阵树定理与方向约定}
实现见第~\pageref{compact-MatrixTree}~页。
返回的是每棵树的边权乘积，再对所有树求和，并对素数取模。
无权计数只需把所有边权设为 1。
重边是不同选择，自环不属于任何生成树；单点图的空树乘积为 1。
零权和负权都按域上的运算处理，因此结果为零可能来自抵消，不能作为连通性判定。

无向边 $uv$ 权值 $w$ 对拉普拉斯矩阵贡献为
$L_{uu},L_{vv}$ 各加 $w$，$L_{uv},L_{vu}$ 各减 $w$。
删除 root 行和列后求行列式，与 root 的选择无关。
对于每个点都能沿有向边走到根的树形图，
边 $u\to v$ 贡献 $L_{uu}\mathrel{+}=w$、$L_{uv}\mathrel{-}=w$；
删除根行列后得到内向树形图的权值和。
外向树形图先反转所有边，再使用同一构造。

例如只有一条边 $0\to1$ 时，内向根 1 的答案为该边权，内向根 0 为零；
外向根 0 的答案为该边权。
\href{https://www.luogu.com.cn/problem/P6178}{P6178} 指定的是以 1 为根的外向树，
驱动把点号减一，并选择 \texttt{away\_from\_root}。
只直接构造 $(n-1)\times(n-1)$ 余子式，时间 $O(n^3+m)$、额外空间 $O(n^2)$。
行列式复用本书的素数模高斯消元，不支持将 mod 随意改成合数。

测试独立枚举 $n-1$ 条边的所有选择，检查连通性或根向约束，
用任意精度整数计算边权乘积之和，再分别对多个素数取模。
两套实现覆盖所有根、三种方向、自环重边、负权及 64 位权值边界，
并通过 220 点完全图的加权 Cayley 公式检查。
P6178 完整驱动另测 300 点、十万条边以及输入输出约定。
普通与 ASan/UBSan 均通过。
两套代码已在 P6178 通过全部 11 个测试点：
\href{https://www.luogu.com.cn/record/297619886}{动态版记录 297619886}、
\href{https://www.luogu.com.cn/record/297620263}{传统版记录 297620263}。
单点最大用时分别为 64 ms、62 ms，内存峰值为 2.85 MB、2.40 MB。
在线测试覆盖无向与根为 1 的外向树；其他根、内向树和负权仍使用本地证据。
定理背景可参考 \href{https://oi-wiki.org/graph/matrix-tree/}{OI Wiki 矩阵树定理}。

